\chapter{NKLogger : dire ce qui se passe}

Un programme graphique n'a pas de console. Quand il se comporte mal, la seule
trace qui reste est celle qu'il a écrite lui-même. Ce chapitre montre comment
l'écrire --- et surtout comment ne pas la rendre fausse, ce qui arrive plus
souvent qu'on ne croit.

\section{Les six niveaux}

\begin{center}
\begin{tabular}{@{}lll@{}}
\toprule
\textbf{Niveau} & \textbf{Pour quoi} & \textbf{Qui le lit} \\
\midrule
\texttt{Trace} & le détail fin d'une boucle & vous, en cherchant \\
\texttt{Debug} & un état intermédiaire & vous, en développant \\
\texttt{Info} & une étape franchie & tout le monde \\
\texttt{Warn} & anormal, mais on continue & tout le monde \\
\texttt{Error} & la fonction a échoué & tout le monde \\
\texttt{Fatal} & on ne peut pas continuer & tout le monde \\
\bottomrule
\end{tabular}
\end{center}

\begin{nkretenir}
Le niveau n'est pas une nuance de politesse : c'est ce qui décide de \emph{ce
qu'on voit} quand on filtre. Un message important écrit en \texttt{Trace}
disparaît le jour où l'on relève le filtre --- et c'est précisément le jour où
l'on cherche quelque chose.
\end{nkretenir}

% Pas de \alerte{} dans un titre : c'est un argument MOBILE (il part vers la
% table des matieres et l'en-tete courant), et le dessin TikZ n'y survit pas --
% « Undefined control sequence \pgf@temp ». Le panneau qui etait ici ne
% s'imprimait de toute facon pas ; un titre de section n'en a pas besoin.
\section{Deux familles de messages, et on ne les mélange pas}

C'est le piège le plus coûteux du module, et il ne plante jamais.

\begin{nkcode}{Les deux familles}
// FAMILLE 1 -- suffixe `f`, style printf. Le `\n` final est EXPLICITE.
logger.Infof("[nkdames] %d coups legaux en %.2f ms\n", n, duree);

// FAMILLE 2 -- sans suffixe, marqueurs INDEXES {i}. Pas de `\n` a ecrire.
logger.Info("[nkdames] {0} coups legaux en {1} ms", n, duree);
\end{nkcode}

\begin{nkpiege}
\textbf{Un marqueur \texttt{\{0\}} passé à \texttt{Infof} n'est pas interprété :
il s'affiche littéralement.} Un \texttt{\%s} passé à \texttt{Info} fait de même.

Dans les deux cas \textbf{rien ne plante, rien n'avertit} : le message sort,
faux, et ressemble à un message. On lit alors une trace qui ment, et l'on cherche
le défaut ailleurs.

Choisissez la famille \emph{avant} d'écrire la chaîne, et n'en changez pas en
cours de ligne.
\end{nkpiege}

\begin{nkpiege}[Le piège des accolades, et il a coûté trois jours]
\textbf{Un texte qui contient des accolades ne passe pas par un formateur à
marqueurs.} La famille 2 prend toute accolade pour un marqueur : un bloc de code,
un JSON, une règle CSS y perdent leur corps --- \emph{en silence}.

\begin{nkcode}{Ce qui arrive reellement}
Entree : "uniform BlocA {\n    vec4 va;\n} uA;"
Sortie : "uniform BlocA 0 uA;"
\end{nkcode}

Et le second membre compte autant que le premier : \textbf{le texte amputé reste
valide} pour l'étape suivante. La panne sort donc à l'autre bout de la chaîne,
des heures plus tard, sous la forme d'une valeur qui n'arrive pas.

Pour ces textes-là : \texttt{Infof} (famille 1), ou une écriture directe. Jamais
la famille 2.
\end{nkpiege}

\section{Où le journal atterrit}

\begin{center}
\begin{tabular}{@{}ll@{}}
\toprule
\textbf{Fichier} & \textbf{Contenu} \\
\midrule
\texttt{logs/app\_<date>\_<heure>\_<pid>.log} & un par lancement, conservé \\
\texttt{logs/app.log} & la course \emph{courante} seule, tronquée au lancement \\
\bottomrule
\end{tabular}
\end{center}

\begin{nkpiege}
\textbf{Une application de ce dépôt n'écrit pas sur la sortie standard.} Elle
écrit par NKLogger dans \texttt{logs/}, \emph{relativement au répertoire
courant} --- jamais à côté du binaire, jamais sur la console en Release.

Conséquence : capturer la sortie standard pour prouver qu'une démo tourne mesure
\textbf{le bon programme sur le mauvais canal}, et fait conclure « démo muette »
sur du code sain. Le diagnostic qui tranche en trois secondes :

\begin{nkcode}{Commande}
find . -name "*.log" -mmin -30
\end{nkcode}
\end{nkpiege}

\begin{nkretenir}
\textbf{Un journal par lancement, et c'est ce qui rend une mesure lisible.}
Avant, \texttt{app.log} accumulait indéfiniment : un fichier réel portait six
lancements sur deux jours, avec vingt-six heures entre deux lignes consécutives
et aucun marqueur. On attribuait à une course les lignes de la précédente.

Les vingt derniers journaux sont conservés, puis purgés au démarrage.
\texttt{NKENTSEU\_LOG\_KEEP} change le nombre ; \texttt{0} désactive la purge.
Un \emph{compte} et non un âge : un âge ne borne rien --- six courses tiennent
dans vingt minutes.
\end{nkretenir}

\section{Ce qu'un bon message contient}

\begin{nkretenir}
\textbf{Un message utile porte un chiffre et sa provenance.}

\begin{center}
\begin{tabular}{@{}ll@{}}
\toprule
\textbf{Inutile} & \textbf{Utile} \\
\midrule
« erreur de chargement » & « [nkimage] hero.png : 0 octet lu sur 4096 attendus » \\
« ça marche » & « [banc] 16 cas sur 16 : VERT » \\
« lent » & « [rendu] 42,3 ms/image, budget 16,6 » \\
\bottomrule
\end{tabular}
\end{center}

Le préfixe entre crochets dit \emph{qui parle}. Sans lui, une trace de cent
lignes ne se démêle plus.
\end{nkretenir}

\begin{nkbilan}
Six niveaux, et le niveau décide de ce qu'on verra en filtrant. \textbf{Deux
familles de messages} --- \texttt{Infof} en \texttt{printf}, \texttt{Info} en
marqueurs \texttt{\{0\}} --- qu'on ne mélange jamais, parce que l'erreur ne
plante pas : elle produit un message faux. Aucun texte à accolades ne passe par
la famille 2. Le journal va dans \texttt{logs/}, un fichier par lancement, jamais
sur la sortie standard. Et un message utile porte un chiffre.
\end{nkbilan}

\begin{nkquestions}
\begin{enumerate}
\item Quelles sont les deux familles de messages, et comment les reconnaît-on ?
\item Que se passe-t-il si l'on passe \texttt{\{0\}} à \texttt{Infof} ?
\item Pourquoi capturer la sortie standard ne montre-t-il rien ?
\item Pourquoi un journal par lancement plutôt qu'un fichier qui s'accumule ?
\item Pourquoi la rétention se compte-t-elle en \emph{nombre} et non en âge ?
\end{enumerate}
\end{nkquestions}

\begin{nkpratique}
Instrumentez une petite application : au démarrage, journalisez la taille de la
fenêtre, le backend graphique retenu et la zone sûre. Toutes les cent images,
journalisez la durée moyenne d'une image.

Puis lancez-la, ouvrez le fichier produit dans \texttt{logs/}, et vérifiez que
vous pouvez répondre à cette question sans relancer : \emph{combien d'images par
seconde a-t-elle tenu ?}
\end{nkpratique}

\begin{nkdemo}
Prouvez que mélanger les familles produit un message faux sans rien casser.

Écrivez la même information deux fois, une fois avec le mauvais formateur :

\begin{nkcode}{A tester}
logger.Infof("valeur = {0}\n", 42);
logger.Info("valeur = %d", 42);
\end{nkcode}

Recopiez les deux lignes telles qu'elles sortent du journal. Puis faites de même
avec un texte contenant une accolade, et montrez ce qui disparaît.
\end{nkdemo}

\begin{nklogique}
Un journal en mode ajout contient six lancements. Une ligne dit
\texttt{[rendu] 42,3 ms/image}.

Quelles informations vous manquent pour savoir si cette ligne concerne la course
que vous venez de lancer ? Énumérez-les.

Puis dites laquelle, à elle seule, suffirait --- et pourquoi un fichier par
lancement les rend toutes inutiles.
\end{nklogique}
